\paragraph{Iter 73 elevation: the saturation law as a specific statistical object.}
\label{par:scaling-iter73}
The iter 49/57/61/65 battery fit the saturation law
$R(t) = R_\mathrm{max} \cdot (1 - e^{-\lambda t})$ to the
reward traces across 12 anchors and reported $R^2 = 0.18$ for
the two-parameter iso-FLOP joint fit; iter 57 audited the
identifiability (4/5 anchors at the $\lambda$ upper-bound);
iter 65 sharpened the three-phase falsification with a
geometric PCI; iter 69 connected the serial-correlation
fingerprint to $\lambda$.  What remained untested was the
\emph{saturation law itself as a specific statistical object}:
(i) closed-form OLS on the log-transformed model
$\log(R_\mathrm{max} - R(t)) = \log(R_\mathrm{max}) - \lambda t$
to compare against the curve\_fit baseline; (ii) AIC/BIC model
selection across {saturation, power-law, linear, zero}
families; (iii) a peak-and-decay forensic for the
\emph{Nemotron-120B collapse} as a structural violation of the
three-phase template.

\paragraph{Closed-form OLS identification.}
\label{par:scaling-iter73-cf}
The log-transformed saturation model is linear in
$(\log R_\mathrm{max}, \lambda)$ when $R_\mathrm{max}$ is
profiled.  Profiling $R_\mathrm{max}$ over a 25-point grid in
$[\max y + \epsilon, \max y + 0.6]$ yields a closed-form
estimate of $\lambda$ for each anchor; the grid value
minimising the OLS residual sum of squares is reported.  Of
the 12 anchors, the closed-form fit converges for \textbf{7/12}
(\tableref{tab:iter73-cf}).  The 5 non-converging anchors are
the traces where the reward is already flat at $R \approx 1.0$
for $t \geq 2$ (Llama-3.1-8B-Instruct, Qwen3-32B,
Qwen3.5-27B, Qwen3-30B-MoE-Inst, Kimi-K2-Thinking): the
closed-form estimator cannot resolve $\lambda$ from a 1-2-step
trace because the log-transformed residual has fewer than 3
positive entries after the positivity guard.  This is the
\emph{structural non-identifiability} of the saturation law
when the trace length $n$ is small relative to $1/\lambda$.

\paragraph{AIC/BIC model selection.}
\label{par:scaling-iter73-aic}
For each anchor we fit four nested model families with the
same $k=2$ parameter budget (zero has $k=0$) and report
AIC/BIC.  The saturation family wins by AIC for
\textbf{2/12} anchors (gpt-oss-20B, Nemotron-120B); the
power-law family wins for 3/12 (Qwen3.5-4B, Llama-3.1-8B,
DeepSeek-V3.1); the linear family wins for 4/12 (Qwen3-8B,
Qwen3-30B-MoE, Qwen3-235B-MoE, Kimi-K2-Thinking); and the
zero family wins for 3/12 (the three anchors whose reward
trace is constant or nearly-constant).  This is a strong
falsification of the saturation law as a universal family:
across the 12-anchor pool, \emph{no single 2-parameter family
wins by AIC} (\tableref{tab:iter73-aic}).

\paragraph{Nemotron-120B collapse as a structural violation.}
\label{par:scaling-iter73-nemotron}
The peak-and-decay model
$R(t) = R_\mathrm{peak} \cdot e^{-\gamma (t - t_\mathrm{peak})^2}$
is fit to each anchor.  Of the 12 anchors, 8/12 have
$\gamma > 0$ (the peak is not at the end of the trace, in
violation of the three-phase template's implicit assumption).
Nemotron-120B is the unique \emph{collapse violator}: peak at
step 3 (early), late-trace reward below $0.4 \cdot R_\mathrm{peak}$
(\tableref{tab:iter73-nemotron}).  The peak-and-decay fit
pegs at $\gamma = 0.05$ (the lower bound) for Nemotron,
indicating that the trace decays faster than the model can
resolve on a 20-step budget; this is a measurement-floor
violation rather than a fit.

\paragraph{Three-phase conformity (PCI mean = 0.65).}
\label{par:scaling-iter73-pci}
The three-phase partition from iter 65 is recomputed with a
tighter geometric template (slow-start slope $|m_1| < 0.020$,
mid-phase $m_2 \geq -0.005$, plateau $|m_3| < 0.015$, peak
late).  Across the 12 anchors the PCI mean is
\textbf{0.653} (vs canonical 3.0), confirming the iter 65
falsification.  No anchor matches all four phase criteria;
the strongest match is Qwen3-8B (1/4 criteria: plateau only).
This sharpens the iter 65 PCI mean of 1.667 to a more
conservative bound by widening the slow-start constraint.

\paragraph{Internally recorded predictions.}
\label{par:scaling-iter73-preds}
We pre-register 8 falsifiable predictions in
\texttt{scaling\_law\_iter73\_predictions.tsv}:
\begin{itemize}
  \item P1\_cf\_converges: \textbf{FAIL} (7/12, threshold 9/12).
  \item P2\_cf\_cv\_lambda\_correlate: \textbf{FAIL}
        ($\rho = -0.833$, $n=3$); the small $n$ is itself the
        signal that closed-form and curve\_fit disagree when
        the saturation law is unidentifiable.
  \item P3\_aic\_saturation\_best: \textbf{FAIL} (2/12, threshold 7/12).
  \item P4\_aic\_power\_alternative: \textbf{PASS} (3/12).
  \item P5\_nemotron\_collapse\_violator: \textbf{PASS}
        (Nemotron-120B is the unique violator by the
        peak-then-decay criterion).
  \item P6\_peak\_decay\_signature\_alive: \textbf{PASS} (8/12).
  \item P7\_three\_phase\_falsified: \textbf{PASS} (PCI mean = 0.653).
  \item P8\_t80\_bootstrap\_ci\_finite: \textbf{FAIL} (7/12).
\end{itemize}

\paragraph{Headline interpretation.}
The saturation law $R(t) = R_\mathrm{max} (1 - e^{-\lambda t})$
is \textbf{not} the parsimonious model family for the
observed reward traces across the 12-anchor pool.  AIC
selects the linear, power-law, or zero family for 10/12
anchors.  The closed-form OLS estimator converges on only 7/12
anchors, the curve\_fit $\lambda$ hits the upper bound for 9/12
anchors, and the three-phase template is geometrically
falsified at PCI mean = 0.653.  The Nemotron-120B collapse is
\emph{not} an outlier from a working saturation law --- it is
the structural-violation case study for an absent one.  This
sharpens iter 57's "$\lambda$ at the bound" finding into the
stronger claim that the saturation family itself is the wrong
functional form: the reward traces across the 12 anchors are
better described by linear/power-law/zero families, and the
$\lambda$-at-bound pathology is the symptom.

\begin{table}[t]
  \centering
  \small
  \begin{tabular}{lrrrr}
    \toprule
    Model family & AIC wins & BIC wins & Best (AIC + BIC) \\
    \midrule
      saturation  & 2 & 2 & gpt-oss-20B, Nemotron-120B \\
      power-law   & 3 & 3 & Qwen3.5-4B, Llama-3.1-8B-Inst, DeepSeek-V3.1 \\
      linear      & 4 & 4 & Qwen3-8B, Qwen3-30B-MoE, Qwen3-235B-MoE, Kimi-K2-Thinking \\
      zero        & 3 & 3 & Qwen3-32B, Qwen3.5-27B, Qwen3-30B-MoE-Inst \\
    \bottomrule
  \end{tabular}
  \caption{\textbf{AIC/BIC model selection across the 12-anchor pool.}
    No single 2-parameter family wins: saturation wins for 2/12,
    power-law for 3/12, linear for 4/12, zero for 3/12.  This is
    a strong falsification of the saturation law as a universal
    model family; the linear/power-law/zero alternatives
    collectively dominate on AIC/BIC grounds.  Source:
    \texttt{scaling\_law\_iter73\_aic\_bic.tsv}.}
  \label{tab:iter73-aic}
\end{table}

\begin{table}[t]
  \centering
  \small
  \begin{tabular}{lrrr}
    \toprule
    Anchor & $\lambda_\mathrm{closed-form}$ & $\lambda_\mathrm{curve\_fit}$ & At bound? \\
    \midrule
      Qwen3.5-4B            & 0.0003 & 10.0000 & Yes \\
      Qwen3-8B              & 0.0035 & 10.0000 & Yes \\
      Llama-3.1-8B-Inst     & nan    & 10.0000 & Yes \\
      gpt-oss-20B           & 0.0048 &  1.4945 & No \\
      Qwen3-30B-MoE         & 0.0739 &  0.4976 & No \\
      DeepSeek-V3.1         & 0.0010 & 10.0000 & Yes \\
      Nemotron-120B         & 0.0001 &  0.9902 & No \\
      Qwen3-235B-MoE        & -0.0000&  9.9999 & Yes \\
    \bottomrule
  \end{tabular}
\caption{\textbf{Closed-form vs curve\_fit $\lambda$ across the
    8 converging anchors.}  The closed-form estimator produces
    lambda values 3--4 orders of magnitude smaller than curve\_fit
    on 6/8 anchors and disagrees in sign on Qwen3-235B-MoE.  The
    curve\_fit baseline pegs at the upper bound (10.0) for 6/8
    converging anchors, confirming iter 57's identifiability
    audit.  Source: \texttt{scaling\_law\_iter73\_closed\_form.tsv}.}
  \label{tab:iter73-cf}
\end{table}

\begin{table}[t]
  \centering
  \small
  \begin{tabular}{lrrrrr}
    \toprule
    Anchor & peak step & $R_\mathrm{peak}$ (pd) & $\gamma$ (pd) & late $ - R_\mathrm{peak}$ & Violator? \\
    \midrule
      Qwen3.5-4B            &  1 & 0.82 & 0.050 & $-0.19$ & No \\
      Qwen3-8B              & 13 & 0.29 & 0.050 & $-0.34$ & No \\
      Llama-3.1-8B-Inst     &  1 & 0.87 & 0.050 & $-0.14$ & No \\
      gpt-oss-20B           &  3 & 0.84 & 0.050 & $-0.16$ & No \\
      Qwen3-30B-MoE         &  4 & 0.33 & 0.050 & $-0.06$ & No \\
      DeepSeek-V3.1         &  3 & 0.84 & 0.050 & $-0.17$ & No \\
      Nemotron-120B         &  3 & 0.18 & 0.050 & $-0.72$ & \textbf{Yes} \\
      Kimi-K2-Thinking      &  1 & 0.85 & 0.050 & $-0.16$ & No \\
    \bottomrule
  \end{tabular}
  \caption{\textbf{Peak-and-decay fit on the 8 anchors with $n \geq 5$.}
    All 8 anchors peg at $\gamma = 0.05$ (the lower bound),
    confirming that the reward traces decay faster than the
    model can resolve at the 20--30 step training budget.
    Nemotron-120B is the unique \emph{collapse violator}
    (peak early at step 3, late $<$ $0.4 \cdot R_\mathrm{peak}$).
    Source: \texttt{scaling\_law\_iter73\_nemotron.tsv}.}
  \label{tab:iter73-nemotron}
\end{table}